% ============================================================
%  Příprava na maturitu – Mechatronika – Okruh 10
% ============================================================

\documentclass[a4paper,10pt]{article}

\usepackage[T1]{fontenc}
\usepackage[utf8]{inputenc}
\usepackage[czech]{babel}
\usepackage[left=2.0cm, right=1.8cm, top=1.8cm, bottom=1.8cm]{geometry}
\usepackage{parskip}
\usepackage{titlesec}
\usepackage{xcolor}
\usepackage{hyperref}
\usepackage{tabularx}
\usepackage{booktabs}
\usepackage{tikz}
\usepackage{mdframed}
\usepackage{amsmath}
\usepackage{enumitem}
\usepackage{multicol}

\usetikzlibrary{arrows.meta, positioning}

% --- Barvy ---
\definecolor{accent}{HTML}{03254E}
\definecolor{lightblue}{HTML}{EAF2F8}
\definecolor{lightyellow}{HTML}{FDFBE4}
\definecolor{lightgreen}{HTML}{EAF7EA}
\definecolor{lightred}{HTML}{FDE8E8}

% --- Boxy ---
\newmdenv[backgroundcolor=lightblue,linecolor=accent!40,roundcorner=4pt,innertopmargin=6pt,innerbottommargin=6pt]{pojmybox}
\newmdenv[backgroundcolor=lightyellow,linecolor=accent!30,roundcorner=4pt,innertopmargin=6pt,innerbottommargin=6pt]{vysvetlenibox}
\newmdenv[backgroundcolor=lightgreen,linecolor=accent!30,roundcorner=4pt,innertopmargin=6pt,innerbottommargin=6pt]{vzorcebox}
\newmdenv[backgroundcolor=lightred,linecolor=accent!30,roundcorner=4pt,innertopmargin=6pt,innerbottommargin=6pt]{durazbox}

\setlist[itemize]{noitemsep, leftmargin=1.4em, topsep=2pt}
\setlist[enumerate]{noitemsep, leftmargin=1.4em, topsep=2pt}

\begin{document}

% ============================================================
\section*{\large Okruh 10 – Průmyslové komunikace: komunikační sběrnice a SCADA systémy}

\begin{pojmybox}
\textbf{Klíčové pojmy:}
signál, A/D a D/A převodník, modulace, Manchester kódování, sběrnice, UART, I2C, SPI, CAN,
RS-232, RS-485, Modbus, Ethernet, SCADA, HMI, master/slave, topologie sítě, simplex/duplex, synchronní/asynchronní komunikace
\end{pojmybox}

\begin{vysvetlenibox}

% ----------------------------------------------------------------
\subsection*{1. Průmyslová komunikace}

Komunikace mezi průmyslovými zařízeními a periferiemi (PLC, senzory, pohony, roboty, HMI, SCADA). Probíhá elektronicky po síti nebo sběrnici a způsob komunikace je standardizován.

\textbf{Požadavky:} otevřenost, spolehlivost, bezpečnost, vrstevnost, deterministická komunikace, odolnost proti rušení, jednoduchá diagnostika.

% ----------------------------------------------------------------
\subsection*{2. Signál}

Fyzikální jev (veličina), která přenáší informaci.

\textbf{Dělení:}
\begin{itemize}
  \item spojitý (analogový) × diskrétní (digitální)
  \item periodický × neperiodický
\end{itemize}
\textbf{Příklady:} optický, akustický, elektrický, mechanický, hydraulický.

% ----------------------------------------------------------------
\subsection*{3. A/D a D/A převodníky}

\begin{vzorcebox}
\textbf{A/D převodník – paralelní:}
Referenční napětí je rozděleno do kvantizačních hladin pomocí odporového děliče. Operační zesilovače porovnají vstupní signál s hladinami a po uplynutí vzorkovací periody odešlou výsledek do dekodéru, který nastaví digitální výstup.

\medskip
\textbf{D/A převodník – váhová struktura:}
U\textsubscript{REF} je pomocí napěťových děličů rozdělen na napěťové úrovně. Řídící člen spíná spínače dle potřebného napětí; výsledné napětí je dáno součtem všech úrovní. Výstupní OZ zajišťuje nízký vstupní odpor (ideálně $R_{vst} = \infty$) a nízký výstupní odpor (ideálně $R_{vyst} = 0$).

\medskip
\textbf{Rozlišení převodníku:}
\[
\Delta U = \frac{U_{REF}}{2^n}
\]
kde $U_{REF}$ = referenční napětí, $n$ = počet bitů. Vyšší $n$ = vyšší přesnost.
\end{vzorcebox}

% ----------------------------------------------------------------
\subsection*{4. Modulace / demodulace}

Metoda, při které se upravují vlastnosti signálu (frekvence, amplituda, fáze), ale nikoliv přenášená informace.

\begin{itemize}
  \item \textbf{AM} – amplitudová modulace
  \item \textbf{FM} – frekvenční modulace
  \item \textbf{PM} – fázová modulace
  \item \textbf{Pulsní} modulace
\end{itemize}
\textbf{Demodulace} – zpětný převod modulovaného signálu na původní. Využití: Wi-Fi, rádio, televizní signál, mobilní sítě.

% ----------------------------------------------------------------
\subsection*{5. Kódování Manchester}

Používá se k přenosu dat přes Ethernet na fyzické vrstvě ISO/OSI. Princip: každý bit původního signálu je „rozpůlen" na dvě části.

\begin{itemize}
  \item Pokud je aktuální bit 1 a následující bude 0 → kódovaný signál přejde do 0 již v polovině bitu 1.
  \item Pokud je aktuální bit 0 a následující bude 1 → kódovaný signál přejde do 1 již v polovině bitu 0.
  \item Když se hodnoty nemění → kódovaný signál kmitá (ale při poslední změně musí nejprve přejít do správné logické hodnoty).
\end{itemize}

% ----------------------------------------------------------------
\subsection*{6. Komunikační sběrnice}

Soustava vodičů přenášejících data stejného charakteru. Řídí ji \textbf{master} (řadič), ostatní prvky jsou \textbf{slave} (podřízené).

\textbf{Parametry sběrnic:}
\begin{itemize}
  \item \textbf{Přenosová rychlost} [bit/s] – max. počet bitů za 1 s
  \item \textbf{Šířka sběrnice} – počet paralelních vodičů
  \item \textbf{Taktovací kmitočet} – základní frekvence (počet hodinových impulsů za 1 s)
  \item \textbf{Délka vedení} – maximální vzdálenost
\end{itemize}

\textbf{Rozdělení sběrnic:}
\begin{itemize}
  \item podle druhu signálů: datové, adresové, řídící, stavové
  \item podle počtu vodičů: sériové, paralelní
\end{itemize}

% ----------------------------------------------------------------
\subsection*{7. Směr a synchronizace přenosu}

\begin{center}
\begin{tabularx}{0.95\linewidth}{@{} l X X @{}}
\toprule
\textbf{Typ} & \textbf{Popis} & \textbf{Příklad} \\
\midrule
Simplex & pouze jedním směrem & ovladač TV \\
Half duplex & obousměrně, ale ne současně & vysílačky \\
Full duplex & obousměrně současně & telefonní hovor, Ethernet \\
\bottomrule
\end{tabularx}
\end{center}

\begin{itemize}
  \item \textbf{Synchronní} komunikace – přenos dat řízen hodinovými impulsy (clock)
  \item \textbf{Asynchronní} komunikace – využívá potvrzující signály (start/stop bity), clock není potřeba
\end{itemize}

% ----------------------------------------------------------------
\subsection*{8. UART}

Jednoduchá asynchronní sériová komunikace. Zařízení musí mít nastavenou stejnou \textbf{baudovou rychlost} (počet bitů za sekundu).

\begin{itemize}
  \item 2 vodiče: \textbf{TX} (vysílání) a \textbf{RX} (příjem)
  \item Rámec: start bit + data (7–8 bitů) + parita + stop bit(y)
  \item \textbf{Parita lichá:} pokud je počet jedniček sudý, přidáme jednu
  \item \textbf{Parita sudá:} pokud je počet jedniček lichý, přidáme jednu
\end{itemize}
\textbf{Použití:} mikrokontroléry, debug konzole, GPS moduly.

% ----------------------------------------------------------------
\subsection*{9. RS-232 vs RS-485}

\begin{center}
\begin{tabularx}{0.97\linewidth}{@{} l X X @{}}
\toprule
\textbf{Vlastnost} & \textbf{RS-232} & \textbf{RS-485} \\
\midrule
Typ přenosu & nesymetrický & diferenciální (vodiče A/B) \\
Počet zařízení & point-to-point (2) & více zařízení na sběrnici \\
Max. vzdálenost & cca 15\,m & až 1200\,m \\
Odolnost proti rušení & nízká & vysoká (EMI) \\
Použití & PC komunikace, sériový port & průmysl, Modbus RTU \\
\bottomrule
\end{tabularx}
\end{center}

RS-485 používá diferenciální přenos (napěťový rozdíl vodičů A a B), díky čemuž je výrazně odolnější vůči elektromagnetickému rušení.

% ----------------------------------------------------------------
\subsection*{10. I²C (Internal Integrated Circuit Bus)}

Synchronní interní datová sběrnice.

\begin{itemize}
  \item 2 vodiče: \textbf{SDA} (data) + \textbf{SCL} (hodinový signál generovaný masterem)
  \item Adresování zařízení – až 128 (7 bit) nebo 1024 (10 bit) čipů na sběrnici
  \item Přenosová rychlost: 100\,kHz / 400\,kHz / 1\,MHz
  \item Nevýhoda: ztráta dat při větší vzdálenosti
  \item \textbf{Start bit:} SDA z 1→0, SCL v logické 1
  \item \textbf{Stop bit:} SDA z 0→1, SCL v logické 1
\end{itemize}
\textbf{Použití:} EEPROM, senzory, LCD displeje, RTC, A/D a D/A převodníky.

% ----------------------------------------------------------------
\subsection*{11. SPI (Serial Peripheral Interface)}

Velmi rychlá synchronní komunikace – 4 vodiče:
\begin{itemize}
  \item \textbf{MOSI} – Master Out Slave In
  \item \textbf{MISO} – Master In Slave Out
  \item \textbf{SCK} – hodinový signál
  \item \textbf{CS/SS} – výběr zařízení (každé slave má vlastní CS vodič)
\end{itemize}
\textbf{Výhody:} vysoká rychlost, jednoduchost. \quad \textbf{Nevýhoda:} každé zařízení potřebuje vlastní CS vodič.\\
\textbf{Použití:} displeje, SD karty, A/D převodníky, flash paměti.

% ----------------------------------------------------------------
\subsection*{12. CAN sběrnice}

Používaná v automobilech a průmyslu.

\begin{itemize}
  \item Diferenciální přenos – vysoká odolnost proti rušení
  \item Prioritizace zpráv (ID zprávy určuje prioritu)
  \item Možnost více master zařízení
\end{itemize}
\textbf{Použití:} ABS, motorové řídicí jednotky, průmyslové stroje.

% ----------------------------------------------------------------
\subsection*{13. Topologie sítí}

\begin{center}
\begin{tabularx}{0.97\linewidth}{@{} l X @{}}
\toprule
\textbf{Topologie} & \textbf{Vlastnosti a chování při výpadku} \\
\midrule
\textbf{Line (liniová)} & Zařízení v linii. Výpadek rozdělí linii na dvě části – jedna zůstane nefunkční. \\
\textbf{Bus (sběrnicová)} & Všechna zařízení na společných vodičích. Výpadek jednoho zařízení neovlivní ostatní. \\
\textbf{Star (hvězdicová)} & Centrální switch/PLC uprostřed. Výpadek nodu (mimo centrum) neovlivní ostatní. \\
\textbf{Ring (kruhová)} & Zařízení v kruhu. Při výpadku lze k sousedovi dostat i z druhé strany. \\
\textbf{Tree (stromová)} & Větve napojené na centrální řízení. Výpadek zařízení vyřadí jen jeho větev. \\
\textbf{Mesh (síťová)} & Více propojených cest. Výpadek lze obejít jinou cestou. \\
\textbf{Fully Connected} & Každé zařízení propojeno s každým. Nejvyšší spolehlivost, ale velký počet vodičů. \\
\bottomrule
\end{tabularx}
\end{center}

% ----------------------------------------------------------------
\subsection*{14. SCADA systémy}

\textbf{SCADA} (Supervisory Control And Data Acquisition) – systém pro dispečerské řízení a sběr dat (továrny, elektrárny, atd.). Obvykle neplní funkci plnohodnotného ovládacího systému.

\textbf{Proč SCADA?}
\begin{itemize}
  \item Snížení počtu odborných pracovníků – linky pracují samostatně, dispečer jen hlídá
  \item Snížení ceny výroby (méně pracovníků, méně dispečerských středisek)
  \item Snadný dozor nad celou výrobou z jednoho počítače
  \item Zpracování až tisíců vstupních dat
  \item Integrace webových technologií a vzdálený přístup přes internet
  \item Snadné ukládání dat a archivace
  \item SCADA komunikuje s PLC pomocí průmyslových sběrnic (Modbus, Profinet, OPC-UA)
  \item Dnes se začínají dostávat i do rodinných domů (nadšenci / smart home)
\end{itemize}

% ----------------------------------------------------------------
\subsection*{15. HMI (Human Machine Interface)}

Rozhraní člověk–stroj – lokální ovládání konkrétního stroje (dotykový panel, tlačítka, displej). \textbf{HMI leží pod SCADA systémem.}

\textbf{Typy HMI displejů:}
\begin{center}
\begin{tabularx}{0.97\linewidth}{@{} l X X @{}}
\toprule
\textbf{Vlastnost} & \textbf{Odporový displej} & \textbf{Kapacitní displej} \\
\midrule
Princip & Vodiče se spojí při stlačení & Změna kapacity při dotyku vodiče \\
Ovládání & Čímkoliv – rukavice, šroubovák & Pouze vodivým předmětem (rukavice nefungují) \\
Životnost & Nižší & Vyšší \\
Využití & Průmysl (nejčastější) & Moderní aplikace \\
\bottomrule
\end{tabularx}
\end{center}

% ----------------------------------------------------------------
\subsection*{16. Rozdíl SCADA vs HMI}

\begin{center}
\begin{tabularx}{0.97\linewidth}{@{} l X X @{}}
\toprule
\textbf{Vlastnost} & \textbf{HMI} & \textbf{SCADA} \\
\midrule
Uživatel & Dělník / operátor stroje & Dispečer / vedoucí provozu \\
Rozsah & Jeden stroj & Celý provoz / závod \\
Funkce & Lokální ovládání & Dohled + sběr dat \\
Archivace & Omezená & Rozsáhlá \\
Vzdálený přístup & Většinou ne & Ano \\
Hierarchie & Podřízeno SCADA & Nadřazeno HMI \\
Vzájemná kontrola & HMI nemůže zasahovat do SCADA & SCADA může ovládat zařízení přes HMI \\
\bottomrule
\end{tabularx}
\end{center}

\end{vysvetlenibox}

\end{document}
